% generated by task5_pinn_silverbox/make_discussion.py
\newcommand{\keptarch}{$64\times64$ tanh}
\newcommand{\keptsettings}{Settings of the kept version: input scaling by the training standard deviations, $64\times64$ tanh hidden layers, linear output scaled by $\sigma_a$; RK4 at $\Delta t$ with sinc-interpolated midpoint input; stage B 500 steps; stage C Adam, learning rate 0.003 with cosine decay, 1000 steps, batch 32 windows, rollout growing from 16 to 128 samples, PINN weight $\lambda_\text{phys}=1$ with the identified coefficients, gradient clip 1; checkpoint at step 1000 (lowest training free-run RMSE).}
\newcommand{\keptdiscussion}{\paragraph{What the sweep shows.}
\begin{itemize}[leftmargin=*,itemsep=1pt]
\item \textbf{The physics loss is what makes the network extrapolate.} On A2\_h2, without it ($\lambda_\text{phys}=0$, run
L\_pw0) the network fits the training record (0.725\,mV) but reaches 7.627\,mV on the test segment; almost all of
that error sits on the 0.205\,V peak, where the network trained on $|y|\le0.174$\,V under-stiffens (force-curve figure). With the
residual to the identified equation the same network gets 0.632\,mV ($\lambda=1$); $\lambda=0.1$ gives 0.773\,mV and
$\lambda=10$ gives 0.758\,mV. Learning $a,b,c,g$ jointly with the network instead of identifying them first
(L\_joint) gives 4.185\,mV: the coefficients drift towards whatever the network does outside the data (in that run the
cubic coefficient ended at $c=4.14\cdot10^{5}$ against the identified $7.47\cdot10^{5}$), so the residual stops carrying information there.
\item \textbf{Training length matters most.} In every architecture the best setting is h5 (1000 steps);
a rollout of 256 (h3) against 128 (h2) helps for some networks and not for others, and RMSprop (h4) is the worst setting for A1, A2, A3, A4. The training free-run error of the kept run was still going down at the end, so more steps would likely lower it further.
\item \textbf{Architecture.} Best test RMSE per architecture: $2\times32$ tanh: 0.481\,mV (A1\_h5); $2\times64$ tanh: 0.422\,mV (A2\_h5); $3\times64$ tanh: 0.447\,mV (A3\_h5); $2\times64$ SiLU: 0.868\,mV (A4\_h5). The tanh networks with 64 units are
better than the 32-unit one at equal settings; SiLU is clearly worse here.
\item \textbf{Comparison with the bare equation.} Eq.~(1) with its four identified coefficients, integrated the same way, scores
0.405\,mV train / 0.393\,mV test; the linear equation ($c=0$) scores 4.961 / 13.546\,mV.
So the cubic term is essential, and on this excerpt the circuit is described by (1) so well that the kept network
(0.422\,mV) matches the identified equation to within
0.03\,mV without beating it.
The network's force is free to depart from (1) where the training data ask for it, and the physics term keeps it on (1) where
there are no data; the physics read-out below shows the force it settled on.
\item \textbf{Train/test gap.} 0.407\,mV vs 0.422\,mV for the kept version: no sign of a
drift on the test segment. The version is chosen by test score, as the task asks, so the reported test number is slightly
optimistic; inside every run the checkpoint was chosen by the training free-run error only.
\end{itemize}}
